\section{Objectives and Scope}
% =========================================================

The simulator has two main purposes. The first is to provide a fast preliminary design tool for a single bell nozzle, from pre-sized propulsion inputs to a usable geometric contour. The second is to optimize the resulting parametric geometry automatically through a genetic algorithm coupled to the same physical model used by the interactive simulator.

The present formulation focuses on:
\begin{enumerate}
    \item geometric definition of a single bell nozzle contour;
    \item use of a throat radius supplied by the engine pre-sizing stage;
    \item calculation of exit radius from a prescribed expansion ratio;
    \item estimation of an ideal expansion ratio for a chosen operating altitude or ambient pressure;
    \item computation of ideal isentropic flow properties along the nozzle;
    \item three-dimensional visualization by revolution of the two-dimensional contour;
    \item coordinate export for CAD modeling and manufacturing preparation;
    \item optimization of expansion ratio, bell-length fraction and initial divergent-wall angle using an effective ambient thrust coefficient.
\end{enumerate}

The simulator is not intended to replace detailed CFD, method-of-characteristics nozzle design, finite-rate chemistry, conjugate thermal analysis, structural analysis or experimental validation. Instead, it provides a fast, interpretable and modifiable preliminary model. Its strongest value is not absolute prediction accuracy, but the ability to connect geometric parameters to performance trends in a controlled and computationally inexpensive way.

% 
